\section{Introduction}\label{sec:intro}

A neutral-atom processor holds its qubits in an array of optical tweezers. Loading is stochastic, so the array
must be sorted after every load and repaired after every loss, and the operations themselves need atoms moved
between sites. Proposal ANRF/ARG/2026/009937~\cite{proposal} plans to do the steering of the transport
tweezers, and the switching of the addressing beams, with silicon-nitride (SiN) photonic integrated circuits
(PICs) instead of acousto-optic deflectors and spatial light modulators. The proposal fixes the working
points. Transport tweezers at 1064.4~nm (a double-magic wavelength for the Cs $45S_{1/2}$ Rydberg state) and
addressing at 894.6~nm (the Cs D1 line) are focused by a meta-objective of numerical aperture $\geq 0.5$ onto a
4~\textmu m lattice. The transport depth is about 1.2~mK, with about 6.6~mW per tweezer at the atoms. The
targets for the circuits are a steering update rate of $\geq 100$~kHz, switch extinction of $\geq 50$~dB at
894.6~nm, a position noise of $\leq 10$~nm rms, and a corrected field of $\pm 100$~\textmu m. The waveguide
platform is the 400~nm stoichiometric LPCVD SiN available to the project.

Whether such circuits can do the job is not obvious, for three reasons. (i) An optical phased array (OPA) steers
all of its tweezers together. Independent motion of many tweezers along a row, as with the tones of an
acousto-optic deflector, needs a different engine. (ii) A chip forms one row. Rows of independent engines must be
stacked and relayed onto a lattice of 4~\textmu m period without the rows disturbing each other. (iii) The
proposal's array is small, while tweezer arrays of 2,088 and 6,100 sites have been reported
\cite{pichard2024,manetsch2024}. An architecture built on chips must therefore also be judged at $10^4$ to
$10^6$ sites. Each question also carries a hard physical limit: the
handoff of atoms between static and moving traps, heating per move, laser power per tweezer, and losses that
compete with repair.

In this work, we simulate the PIC elements and the system around them, from waveguide modes to the steady
state of a $10^6$-site array. The models are scalar or vector wave optics, beam propagation, 2D FDTD, ray
tracing, classical and quantum transport, and stochastic scheduling of moves and losses. Every number carries
its provenance. Values from the proposal are marked USER, values from sources read during the work SOURCE,
simulation outputs SIM, calculations DERIVED, stated assumptions CHOICE, and values recalled but not verified
RECALLED. All simulated values are valid in the regime simulated and should be reviewed before they enter the
proposal. The code, results and module-by-module documents are in the repository. Volume~II compiles those
documents without rewriting, and \cref{app:repro} maps every result to its script.

The report is organised by physical question. \Cref{sec:platform} covers the waveguide platform and the passive
and active components. \Cref{sec:engine} develops the steering engine for independent tweezers.
\Cref{sec:rows} stacks engines into rows and treats their mutual independence. \Cref{sec:moretweezers} asks how
far the tweezer count per engine and per row can grow. \Cref{sec:transport} simulates transport and handoff
in three dimensions, and \cref{sec:rearrange} schedules rearrangement and repair. \Cref{sec:megascale} scales
the architecture to $10^4$, $10^5$ and $10^6$ sites. \Cref{sec:combine} combines chip stacks into larger
blocks and replaces the ideal collimator by the real one, which exposes a limit of the cylindrical collimator
for tweezers away from the row centre. \Cref{sec:discussion} separates what was shown from what is
model-dependent, and states the main limitations and next steps.
